\section{Introduction}
\label{sec:introduction}

Serving a transformer at a 128k-token context is no longer an arithmetic
problem. On the accelerators most laboratories actually own, the attention
kernel spends the majority of its wall-clock time moving key and value tiles
between high-bandwidth memory and the on-chip scratchpad. Doubling the context
doubles that traffic, and the arithmetic units idle while it drains.

Sparsity is the obvious response, and the literature offers a great deal of
it~\cite{child2019sparse,beltagy2020longformer,zaheer2020bigbird}. Most of that
work fixes a sparsity pattern in advance: a sliding window, a dilated stride, a
handful of global tokens. Fixed patterns are easy to schedule, but they commit
to a decision before the model has seen the prompt, and the heads that matter
for a particular prompt are not known in advance.

We take the opposite position. \sysname{} keeps the attention matrix dense in
principle and decides, per block and per prompt, which blocks can be skipped
without changing the output beyond a stated tolerance. The decision is made by a
rank-4 probe over the query and key tiles, which costs roughly 1.5 percent of
the dense kernel, and it is checked at runtime by a guard that restores dense
computation whenever the probe's residual bound is violated.

\paragraph{Contributions.}
This paper makes three contributions.
\begin{itemize}
  \item A block-importance predictor that estimates the contribution of an
        attention block from a low-rank sketch of its query and key tiles,
        with an error bound that holds for any softmax temperature
        (Section~\ref{sec:design}).
  \item A scheduler that turns those predictions into a memory plan, so that
        skipped blocks are never fetched rather than being fetched and
        discarded (Section~\ref{sec:design}).
  \item An evaluation on three open-weight models at contexts from 8k to 128k
        tokens, including the failure cases where the approximation degrades
        (Section~\ref{sec:evaluation}).
\end{itemize}
